% !PRINTED 166
\weberpage{166}

%<<CONT
two functions $UV$ can be divisible by $v$ only if one of the two
factors is divisible by $v$. The same holds for a product of
several functions, and so from I. there follows the theorem:

\gesperrt{II. A product of several integral rational functions is divisible by an indecomposable function $v$ only if at least one of the factors of the product is divisible by $v$.}

If an integral rational function $U$ is decomposed in two
ways into indecomposable factors,
\[
U = v\,v'\,v''\dots = w\,w'\,w''\dots,
\]
then by II. at least one of the factors $v, v', v''\dots$ must be
divisible by $w$, say $v$. But then, since $v$ too is indecomposable,
$v$ can differ from $w$ only by a constant factor.

Accordingly, if $c$ is this constant factor,
\[
c\,v'\,v''\cdots = w'\,w''\dots,
\]
from which it follows that one of the functions $v', v''\dots$, say $v'$,
is divisible by $w'$, and thus differs from $w'$ only by a constant
factor.

We therefore obtain, as a second consequence of Theorem I:

\gesperrt{III. An integral rational function can, apart from constant factors, be decomposed into indecomposable factors in only one way.}

From this arises the concept of the \gesperrt{greatest common
divisor} of two or more integral rational functions
$U, V\dots$ By this one understands the product of all
indecomposable factors that occur in the decompositions of each of the
functions $U, V\dots$, or the function of highest possible
degree that divides all the functions $U, V\dots$ By III. this
function is, apart from a constant factor, completely determined
for every system of functions $U, V\dots$

Several functions $U, V, W\dots$ are called relatively prime if
no two of them have a common divisor.

All these definitions and theorems are exactly analogous to
very well-known theorems of elementary number theory, which there
appear as consequences of the algorithm of the greatest common
%CONT>>
